\section{幺半群上的全代数}

记号约定:$A$是交换环,$S$是幺半群(乘法记号).

\begin{definition}{局部有限乘法条件}{}
	若对任一$s\in S$,集合$\left\{\left(t,u\right)\in S\times S\middle|tu=s\right\}$有限,则称$S$满足局部有限乘法条件.
\end{definition}

\begin{definition}{全幺半群代数}{}
	设$S$满足局部有限乘法条件.定义$A^{S}$上的乘法为$\left(\left(\alpha_{s}\right)_{s\in S},\left(\beta_{s}\right)_{s\in S}\right)\mapsto\left(\sum\limits_{\substack{t,u\in S\\tu=s}}\alpha_{t}\beta_{u}\right)_{s\in S}$,则$A^{S}$是$A$-含幺结合代数.我们称之为$S$在$A$上的全幺半群代数.
\end{definition}

\begin{remark}{记号和术语说明}{}
	\begin{enumerate}
		\item 幺半群代数$A^{\left(S\right)}$是$A^{S}$的子代数.我们称$A^{\left(S\right)}$是$S$在$A$上的限制幺半群代数.
		\item 记$\left(e_{s}\right)_{s\in S}$是$A^{\left(S\right)}$的典范基.$A^{S}$的任意元素$\left(\xi_{s}\right)_{s\in S}$可以在形式上写作$\sum\limits_{s\in S}\xi_{s}e_{s}$(甚至$\sum\limits_{s\in S}\xi_{s}s$).需要注意的是,这里的求和好并不是代数中的无穷次加法计算,仅仅是\textbf{形式上的记号}.
	\end{enumerate}
\end{remark}

\begin{proposition}{}{}
	若$S$是交换幺半群,则$A^{S}$是$A$-含幺交换结合代数.
\end{proposition}

\begin{proposition}{}{}
	若$T$是$S$的子幺半群,则$A^{T}$典范等同于$A^{S}$的一个子代数.
\end{proposition}

\begin{proposition}{函子性}{}
	设$B$是交换环,$\rho\in\Hom_{\mathsf{Ring}}\left(A,B\right)$,则典范映射$A^{S}\to B^{S},\left(\alpha_{s}\right)_{s\in S}\mapsto\left(\rho\left(\alpha_{s}\right)\right)_{s\in S}$是$A$-代数同态.
\end{proposition}











































































